\documentclass{article}
\usepackage[T1]{fontenc}
\usepackage{lmodern}
\usepackage{graphicx}
\usepackage{amsmath,amssymb}
\usepackage[a4paper,margin=2cm]{geometry}
\usepackage[absolute,overlay]{textpos}
\setlength{\TPHorizModule}{1mm}
\setlength{\TPVertModule}{1mm}
\usepackage[indonesian]{babel}
\usepackage{hyperref}
\setlength{\emergencystretch}{2em}
% unit: Halaman 4

\begin{document}

% === Halaman 4 (llm) ===

\begin{itemize}
    \item $E$ dan $D$ adalah fungsi enkripsi dan dekripsi dari suatu \textit{cipher} blok dengan parameter kunci $K$:
    \[ C = E_K(P) \quad \text{dan} \quad P = D_K(C) \]
    \item Dimisalkan contoh sebuah fungsi $E$ yang sederhana, tapi tidak cukup aman, sbb:
    \begin{enumerate}
        \item XOR-kan blok plainteks $P$ dengan $K$
        \item lalu geser secara \textit{wrapping} bit-bit dari hasil langkah 1 satu posisi ke kiri
    \end{enumerate}
    \[ C = E_K(P) = (P \oplus K) \ll 1 \]
    \item $D$ adalah kebalikan (\textit{invers}) dari $E$. Contoh fungsi $D$ yang sederhana adalah sbb:
    \begin{enumerate}
        \item Geser secara \textit{wrapping} bit-bit ciphertes $C$ satu posisi ke kanan
        \item Lalu XOR-kan blok hasil Langkah 1 dengan $K$
    \end{enumerate}
    \[ P = D_K(C) = (C \gg 1) \oplus K \]
\end{itemize}

\end{document}
